\section{Operacionalización de la Hipótesis}\label{cap:operacionalizacion}

La Tabla~\ref{tab:v4_4} traduce las variables en dimensiones, indicadores e instrumentos medibles. La variable independiente corresponde a la implementación del sistema; la dependiente, a la eficiencia y exactitud de la gestión administrativa.


\vspace{8pt}{\footnotesize\setstretch{1.0}
\begin{longtable}{>{\raggedright\arraybackslash}p{0.1434\dimexpr\linewidth-12\tabcolsep\relax}>{\raggedright\arraybackslash}p{0.1992\dimexpr\linewidth-12\tabcolsep\relax}>{\raggedright\arraybackslash}p{0.1713\dimexpr\linewidth-12\tabcolsep\relax}>{\raggedright\arraybackslash}p{0.1355\dimexpr\linewidth-12\tabcolsep\relax}>{\raggedright\arraybackslash}p{0.1952\dimexpr\linewidth-12\tabcolsep\relax}>{\raggedright\arraybackslash}p{0.1554\dimexpr\linewidth-12\tabcolsep\relax}}
\caption{Operacionalización de las variables}\label{tab:v4_4}\\
\toprule
\textbf{Variable} & \textbf{Definición conceptual} & \textbf{Definición operacional} & \textbf{Dimensión} & \textbf{Indicadores} & \textbf{Escala / instrumento} \\
\midrule\endfirsthead
\toprule
\textbf{Variable} & \textbf{Definición conceptual} & \textbf{Definición operacional} & \textbf{Dimensión} & \textbf{Indicadores} & \textbf{Escala / instrumento} \\
\midrule\endhead
\bottomrule\endfoot
\bottomrule\endlastfoot
\textbf{Independiente (causa):} implementación del sistema web móvil Lush Beauty & Sistema de información que recopila, procesa y distribuye datos para apoyar el control de la organización (Laudon \& Laudon, 2016), cuya calidad en uso se evalúa según ISO/IEC (2011). & Sistema operando en el salón, con registro de servicios contra catálogo, caja diaria y liquidación automática. & 1. Cobertura funcional\newline{}2. Usabilidad percibida & \% de servicios registrados vinculados al catálogo; n.º de funciones operativas del MVP; puntaje SUS (0–100). & Razón / ficha de verificación funcional; ordinal-intervalo / cuestionario SUS (Brooke, 1996). \\
\textbf{Dependiente (efecto):} eficiencia y exactitud de la gestión administrativa & Grado en que los procesos de registro y liquidación se completan con el menor consumo de tiempo y sin desviaciones respecto del resultado correcto (ISO/IEC, 2011; DeLone \& McLean, 2003). & Diferencia entre la línea base retrospectiva y las semanas piloto en tiempo y en discrepancias de cálculo. & 1. Eficiencia temporal\newline{}2. Exactitud del cálculo\newline{}3. Trazabilidad & Minutos por liquidación semanal; segundos por registro de servicio;~\% de líneas con discrepancia de comisión;~\% de liquidaciones con diferencia;~\% de registros con método de pago y bitácora. & Razón / ficha de observación cronometrada y ficha de conciliación contra cálculo de referencia. \\
\end{longtable}}
\nota{Elaboración propia. SUS: System Usability Scale.}

\subsection{Variables e indicadores del soporte para decisiones}

\vspace{8pt}{\small\setstretch{1.0}
\begin{longtable}{>{\raggedright\arraybackslash}p{0.2500\dimexpr\linewidth-8\tabcolsep\relax}>{\raggedright\arraybackslash}p{0.2500\dimexpr\linewidth-8\tabcolsep\relax}>{\raggedright\arraybackslash}p{0.2500\dimexpr\linewidth-8\tabcolsep\relax}>{\raggedright\arraybackslash}p{0.2500\dimexpr\linewidth-8\tabcolsep\relax}}
\caption{Operacionalización ampliada del soporte para decisiones}\label{tab:dss_17}\\
\toprule
\textbf{Variable} & \textbf{Dimensión} & \textbf{Indicador} & \textbf{Medición} \\
\midrule\endfirsthead
\toprule
\textbf{Variable} & \textbf{Dimensión} & \textbf{Indicador} & \textbf{Medición} \\
\midrule\endhead
\bottomrule\endfoot
\bottomrule\endlastfoot
Independiente: uso DSS & Cobertura del proceso & \% de decisiones con alternativas y criterios completos & Lista de verificación \\
Dependiente: calidad del proceso & Eficiencia & Minutos por tarea de comparación & Observación cronometrada \\
Dependiente: calidad del proceso & Trazabilidad & \% decisiones con fuentes, supuestos y aprobación & Auditoría \\
Dependiente: calidad del proceso & Comprensión & Explicación correcta de las razones de la recomendación & Guía breve de evaluación \\
Técnica: motor de escenarios & Exactitud & \% de cálculos coincidentes con referencia independiente & Casos unitarios \\
Técnica: recomendaciones & Seguridad & N.º de acciones financieras efectuadas sin aprobación & Pruebas de permisos \\
Seguimiento & Contraste & Diferencia entre indicador esperado y observado & Registro de decisiones \\
\end{longtable}}
\nota{Elaboración propia. Indicadores de evaluación propuestos, no medidos.}

\subsection{Precauciones de interpretación}

No se inferirá causalidad fuerte sobre mejoras económicas a partir de un solo salón y pocas participantes. Se priorizarán validaciones deterministas del motor matemático y pruebas de tareas repetibles con tiempo y observación. Las encuestas de usabilidad y percepción se presentarán como evidencia descriptiva. Para estudiar impacto comercial real de una promoción se requeriría diseño adicional, autorización y observaciones suficientes.

